\documentclass[12pt]{article}

% Packages
\usepackage[margin=1in]{geometry}
\usepackage{amsmath}
\usepackage{amssymb}
\usepackage{amsthm}
\usepackage{enumitem}

% Source - https://tex.stackexchange.com/a/21229
% Posted by Carsten Thiel
% Retrieved 2026-02-01, License - CC BY-SA 3.0

\theoremstyle{definition}
\newtheorem{exmp}{Example}


% Title information
\title{Mathematics 629\\
Spring 2026\\
Homework Assignment 4}
\author{Alexander Kim}
\date{Due: 2/27/2026}

\begin{document}

\maketitle

\begin{enumerate}

	\item First some terminology: For a series $\sum_{j=0}^{\infty} a_j$ we denote by $s_n$ the partial sums $s_n = \sum_{j=0}^{n} a_j$. The series is \textbf{Ces\`{a}ro summable} if $\lim_{N \to \infty} \frac{1}{N+1} \sum_{n=0}^{N} s_n$ exists. The series is \textbf{Abel summable} if $\lim_{r \to 1^{-}} \sum_{j=0}^{\infty} a_j r^j$ exists in $\mathbb{C}$.

	      Now let $(X, \mathcal{M})$ be a measurable space and $f_j : X \to \mathbb{C}$, $j \in \mathbb{N}_0$, define a sequence of measurable functions. Let
	      \[
		      E_{\text{Ces}} = \left\{ x \in X : \sum_{j=0}^{\infty} f_j(x) \text{ is Ces\`{a}ro summable} \right\}
	      \]
	      \[
		      E_{\text{Abel}} = \left\{ x \in X : \sum_{j=0}^{\infty} f_j(x) \text{ is Abel summable} \right\}
	      \]
	      Show that $E_{\text{Ces}}$ and $E_{\text{Abel}}$ belong to $\mathcal{M}$.

	      \begin{proof}
		      % Your proof here.
	      \end{proof}

	\item Let $f_n : [0,1] \to \mathbb{R}$ be a sequence of measurable functions such that $|f_n(x)| < \infty$ almost everywhere. Show that there is a sequence $c_n$ of positive real numbers such that
	      \[
		      \lim_{n \to \infty} \frac{f_n(x)}{c_n} = 0 \quad \text{for almost every } x.
	      \]

	      \textit{Hint: To start show first that one can pick $c_n$ so that $m(\{x : |f_n(x)/c_n| > 1/n\}) < 2^{-n}$.}

	      \begin{proof}
		      For each $f_n$, let $S_n(c) = \{x \in [0,1]; \vert f_n(x) \vert > \frac{c}{n}\}$. Then
		      as $c \rightarrow \infty$, $m(\cap_{k=1}^n S_k) = 0$ where 
		      $S_k = \{x: \vert f_k(x)\vert  = \infty$ \} since we're assuming all $\vert f_n (x) \vert < \infty$ for all $f_n$. Thus, there exists sufficiently large $c_n$ such that 
		      $$m(E_n = \{x: \frac{\vert f_n(x)\vert}{c_n} > \frac{1}{n} \}) < 2^{-n} < \infty.$$
		      By the Borel-Cantelli Lemma, $\mu(\text{limsup}E_n) = 0$, so for all $n$ greater than
		      some $N$, the set $E_n$ becomes sufficiently small to show $\lim_{n\rightarrow \infty}\frac{f_n(x)}{c_n} = 0$ for almost every $x$.
	      \end{proof}

	\item Let $E \subset \mathbb{R}$ be a compact set and let $d(x, E) = \inf\{|x - y| : y \in E\}$.

	      \begin{enumerate}[label=(\roman*)]
		      \item Show that $x \mapsto d(x, E)$ is a continuous function.

		            \begin{proof}
				    Let $a \in E$ and $x \in E$ such that $\lim \vert x - a \vert < \epsilon$ for all $\epsilon, \delta \in \mathbb{R}$ where $\epsilon = \delta$. Then for any $y \in E$, by the triangle inequality
				    $$ \vert x - y \vert \leq \vert x - a \vert + \vert a - y \vert $$
				    or equivalently
				    $$ d(x, E) \leq \vert x - a\vert + d(a, E).$$
				Similarly, $d(a, E) \leq \vert x - a \vert + d(x, E)$ by the construction of $d(x, E)$. 
				Thus, $\vert d(x,E) - d(a, E) \vert \leq \vert x - a \vert < \epsilon = \delta$ so $x\mapsto d(x, E)$ is continuous by definition.

		            \end{proof}

		      \item Let $m$ be Lebesgue measure on $\mathbb{R}$. Let $O_n = \{x : d(x, E) < 2^{-n}\}$. Show that $\lim_{n \to \infty} m(O_n) = m(E)$.

		            \begin{proof}
			            % Your proof here.
		            \end{proof}
	      \end{enumerate}

\end{enumerate}

\end{document}
